\chapter[Lezione 22 --- 20 novembre 2026]{Lezione 22: derivabilità e derivate delle funzioni fondamentali}
\label{chap:lezione-22}

\section{Rapporto incrementale e derivabilità}
\label{sec:lez22-derivabilita}

Sia \(f\colon I\longrightarrow\R\), con \(I\subseteq\R\)
intervallo, e sia \(x_0\) un punto interno a \(I\)
(ossia esiste \(\delta>0\) tale che
\(D(x_0,\delta)\subseteq I\)).
Per ogni \(x\in I\setminus\{x_0\}\), si considera
la variazione della funzione rispetto alla variazione
della variabile indipendente.

\begin{definition}[Rapporto incrementale]
\label{def:lez22-rapporto-incrementale}
Il \emph{rapporto incrementale} di \(f\) in \(x_0\)
è l'espressione
\[
  \frac{f(x)-f(x_0)}{x-x_0},
  \qquad x\in I\setminus\{x_0\}.
\]
\end{definition}

Il rapporto incrementale è il coefficiente angolare della
retta secante al grafico di \(f\) nei punti
\((x_0,f(x_0))\) e \((x,f(x))\).

\begin{figure}[htbp]
  \centering
  \begin{tikzpicture}[unipd/diagram,x=1.36cm,y=1.10cm]
  \pgfmathsetmacro{\xzero}{0.70}
  \pgfmathsetmacro{\xval}{2.45}
  \pgfmathsetmacro{\yzero}{.21*\xzero*\xzero+.25*\xzero+.85}
  \pgfmathsetmacro{\yval}{.21*\xval*\xval+.25*\xval+.85}
  \pgfmathsetmacro{\msec}{(\yval-\yzero)/(\xval-\xzero)}
  \pgfmathsetmacro{\mtan}{.42*\xzero+.25}

  % Assi
  \draw[->,draw=unipd-muted] (-.46,0)--(3.22,0)
    node[right] {$x$};
  \draw[->,draw=unipd-muted] (0,-.18)--(0,3.58)
    node[above] {$y$};

  % Funzione
  \draw[very thick,draw=unipd-accent,domain=-.30:2.89,samples=100]
    plot (\x,{.21*\x*\x+.25*\x+.85});

  % Secante
  \draw[thick,draw=unipd-primary]
    (-.08,{\yzero+\msec*(-.08-\xzero)})--
    (2.93,{\yzero+\msec*(2.93-\xzero)});

  % Tangente
  \draw[thin,draw=unipd-muted]
    (-.05,{\yzero+\mtan*(-.05-\xzero)})--
    (2.94,{\yzero+\mtan*(2.94-\xzero)});

  % Proiezioni sui due assi
  \draw[densely dashed,draw=unipd-muted,opacity=.65]
    (\xzero,0)--(\xzero,\yzero)
    (\xval,0)--(\xval,\yval)
    (0,\yzero)--(\xzero,\yzero)
    (0,\yval)--(\xval,\yval);

  % Punti
  \fill[fill=unipd-accent] (\xzero,\yzero) circle (2.1pt);
  \fill[fill=unipd-accent] (\xval,\yval) circle (2.1pt);

  % Etichette sugli assi
  \node[below=4pt] at (\xzero,0) {$x_0$};
  \node[below=4pt] at (\xval,0) {$x$};
  \node[left=5pt] at (0,\yzero) {$f(x_0)$};
  \node[left=5pt] at (0,\yval) {$f(x)$};

  % Legenda esterna
  \draw[very thick,draw=unipd-accent]
    (3.30,3.44)--(3.59,3.44);
  \node[anchor=west,font=\sffamily\footnotesize,text=unipd-accent]
    at (3.70,3.44) {funzione};

  \draw[thick,draw=unipd-primary]
    (3.30,3.04)--(3.59,3.04);
  \node[anchor=west,font=\sffamily\footnotesize,text=unipd-primary]
    at (3.70,3.04) {secante};

  \draw[thin,draw=unipd-muted]
    (3.30,2.64)--(3.59,2.64);
  \node[anchor=west,font=\sffamily\footnotesize,text=unipd-muted]
    at (3.70,2.64) {tangente};
\end{tikzpicture}
  \caption{Il rapporto incrementale rappresenta la pendenza
  della secante. Al tendere di \(x\) a \(x_0\), se il limite
  è finito, la secante tende alla retta tangente.}
  \label{fig:lez22-rapporto-incrementale}
\end{figure}

\begin{definition}[Derivabilità in un punto]
\label{def:lez22-derivabilita}
La funzione \(f\) si dice \emph{derivabile in \(x_0\)}
se esiste ed è finito il limite
\[
  f'(x_0):=\lim_{x\to x_0}
    \frac{f(x)-f(x_0)}{x-x_0}\in\R.
\]
Il numero \(f'(x_0)\) è detto \emph{derivata prima}
di \(f\) in \(x_0\).
\end{definition}

Equivalentemente, posto \(h:=x-x_0\), si ha
\[
  f'(x_0)=\lim_{h\to0}\frac{f(x_0+h)-f(x_0)}h.
\]

Se \(f\) è derivabile in \(x_0\), il coefficiente
angolare della retta tangente al grafico in
\((x_0,f(x_0))\) è \(f'(x_0)\); la sua equazione è
\[
  y-f(x_0)=f'(x_0)(x-x_0).
\]
Un rapporto incrementale che tenda a \(+\infty\)
o \(-\infty\) non definisce una derivata finita:
può indicare una tangente verticale, ma non la
derivabilità nel senso della \cref{def:lez22-derivabilita}.

\section{Derivabilità e continuità}
\label{sec:lez22-derivabilita-continuita}

\begin{theorem}[La derivabilità implica la continuità]
\label{thm:lez22-derivabilita-continuita}
Sia \(f\colon I\longrightarrow\R\), con \(I\)
intervallo, e sia \(x_0\) interno a \(I\).
Se \(f\) è derivabile in \(x_0\), allora è continua
in \(x_0\).
\end{theorem}

\begin{proof}
Per ogni \(x\ne x_0\), si ha
\[
  f(x)-f(x_0)
  =\frac{f(x)-f(x_0)}{x-x_0}(x-x_0).
\]
Per la \cref{def:lez22-derivabilita}, il primo fattore
tende a \(f'(x_0)\in\R\), mentre il secondo tende
a \(0\). Dal \cref{thm:lez14-limite-prodotto} segue
\[
  \lim_{x\to x_0}\bigl(f(x)-f(x_0)\bigr)
  =f'(x_0)\cdot0=0.
\]
Pertanto \(\lim_{x\to x_0}f(x)=f(x_0)\) e, per
la \cref{def:lez21-continuita}, \(f\) è continua
in \(x_0\).
\end{proof}

Il viceversa è falso, come mostra il seguente esempio.

\begin{example}[Continuità senza derivabilità]
\label{ex:lez22-valore-assoluto}
La funzione \(f(x)=|x|\) è continua su \(\R\),
ma non è derivabile in \(0\). Infatti,
\[
  \lim_{x\to0^+}\frac{|x|-|0|}{x-0}=1,
  \qquad
  \lim_{x\to0^-}\frac{|x|-|0|}{x-0}=-1.
\]
Poiché i limiti unilateri del rapporto incrementale
sono diversi, la derivata in \(0\) non esiste.
\end{example}

\section{Punti di non derivabilità}
\label{sec:lez22-punti-non-derivabilita}

I limiti unilateri del rapporto incrementale
permettono di distinguere alcuni comportamenti
nei punti in cui la derivata non esiste.

\begin{definition}[Punto angoloso]
\label{def:lez22-punto-angoloso}
Sia \(f\colon I\longrightarrow\R\), con \(x_0\)
interno a \(I\), e sia \(f\) continua in \(x_0\). Il punto \(x_0\) si dice
\emph{punto angoloso} se i due limiti unilateri
\[
  \lim_{x\to x_0^-}\frac{f(x)-f(x_0)}{x-x_0},
  \qquad
  \lim_{x\to x_0^+}\frac{f(x)-f(x_0)}{x-x_0}
\]
esistono finiti, ma sono diversi.
\end{definition}

\begin{definition}[Punto cuspidale]
\label{def:lez22-punto-cuspidale}
Sia \(f\colon I\longrightarrow\R\), con \(x_0\)
interno a \(I\), e sia \(f\) continua in \(x_0\). Il punto \(x_0\) si dice
\emph{punto cuspidale} se i due limiti unilateri
\[
  \lim_{x\to x_0^-}\frac{f(x)-f(x_0)}{x-x_0},
  \qquad
  \lim_{x\to x_0^+}\frac{f(x)-f(x_0)}{x-x_0}
\]
sono infiniti di segno opposto.
\end{definition}

In entrambi i casi la funzione non è derivabile in
\(x_0\), poiché non esiste un limite bilatero finito.
Anche un rapporto incrementale con un solo limite
unilaterale infinito esclude la derivabilità, senza
necessariamente individuare una cuspide.

\clearpage

\begin{figure}[htbp]
  \centering
  \begin{tikzpicture}[unipd/diagram,x=1cm,y=.95cm]
  % Stesse dimensioni e assi per entrambi i grafici.
  \begin{scope}[shift={(0,0)}]
    \draw[->,unipd-muted] (-1.65,0)--(1.65,0) node[right] {$x$};
    \draw[->,unipd-muted] (0,-.18)--(0,2.05) node[above] {$y$};
    \draw[very thick,unipd-primary]
      (-1.35,1.35)--(0,0)--(1.35,1.35);
    \fill[unipd-primary] (0,0) circle (1.7pt);
    \node[below left] at (0,-.12) {$0$};
    \node[align=center,font=\sffamily\footnotesize]
      at (0,-.95) {Punto angoloso\\$f(x)=|x|$};
  \end{scope}
  \begin{scope}[shift={(4.70,0)}]
    \draw[->,unipd-muted] (-1.65,0)--(1.65,0) node[right] {$x$};
    \draw[->,unipd-muted] (0,-.18)--(0,2.05) node[above] {$y$};
    \draw[very thick,unipd-primary,domain=0:1.5,samples=80]
      plot (\x,{1.3*sqrt(\x)});
    \draw[very thick,unipd-primary,domain=-1.5:0,samples=80]
      plot (\x,{1.3*sqrt(-\x)});
    \fill[unipd-primary] (0,0) circle (1.7pt);
    \node[below left] at (0,-.12) {$0$};
    \node[align=center,font=\sffamily\footnotesize]
      at (0,-.95) {Punto cuspidale\\$g(x)=\sqrt{|x|}$};
  \end{scope}
\end{tikzpicture}

  \caption{Punto angoloso per \(f(x)=|x|\) (a sinistra)
  e punto cuspidale per \(g(x)=\sqrt{|x|}\) (a destra),
  entrambi nell'origine.}
  \label{fig:lez22-punti-non-derivabilita}
\end{figure}

\begin{example}[Un punto cuspidale]
\label{ex:lez22-radice-valore-assoluto}
Si consideri \(g(x)=\sqrt{|x|}\), continua in \(0\).
Poiché \(g(0)=0\), si ha
\[
\begin{aligned}
  \lim_{x\to0^+}\frac{g(x)-g(0)}{x}
    &=\lim_{x\to0^+}\frac{1}{\sqrt{x}}=+\infty,\\
  \lim_{x\to0^-}\frac{g(x)-g(0)}{x}
    &=-\lim_{x\to0^-}\frac{1}{\sqrt{-x}}=-\infty.
\end{aligned}
\]
Per la \cref{def:lez22-punto-cuspidale},
\(0\) è un punto cuspidale di \(g\).
\end{example}

\section{Insiemi di continuità e derivabilità}
\label{sec:lez22-insiemi-derivabilita}

\begin{definition}[Insiemi di continuità e derivabilità]
\label{def:lez22-insiemi-continuita-derivabilita}
Sia \(f\colon I\longrightarrow\R\), con \(I\)
intervallo. Si definiscono
\[
  D_{\mathrm{cont}}(f)
  :=\{x\in I:f\text{ è continua in }x\},
\]
\[
  D_{\mathrm{der}}(f)
  :=\{x\in\operatorname{int}(I):f\text{ è derivabile in }x\}.
\]
\end{definition}

Per il \cref{thm:lez22-derivabilita-continuita}, vale
\[
  D_{\mathrm{der}}(f)\subseteq D_{\mathrm{cont}}(f)
  \subseteq I.
\]
L'inclusione inversa non vale in generale, come
mostra l'\cref{ex:lez22-valore-assoluto}.

\section{Funzione derivata prima}
\label{sec:lez22-funzione-derivata}

\begin{definition}[Funzione derivata prima]
\label{def:lez22-funzione-derivata}
Sia \(f\colon I\longrightarrow\R\). La
\emph{funzione derivata prima} di \(f\) è
l'applicazione
\[
  f'\colon D_{\mathrm{der}}(f)\longrightarrow\R,
  \qquad x\longmapsto f'(x).
\]
La derivata si indica anche con \(f^{(1)}\) oppure
con \(\frac{d}{dx}f\).
\end{definition}

Se \(f\) è derivabile in ogni punto interno a \(I\),
allora \(D_{\mathrm{der}}(f)=\operatorname{int}(I)\).
In particolare, se \(I\) è aperto, la derivata è
definita sull'intero intervallo \(I\).

\section{Derivate delle funzioni fondamentali}
\label{sec:lez22-derivate-fondamentali}

I limiti notevoli delle
\cref{sec:lez17-limiti-notevoli,sec:lez19-limite-potenze,sec:lez20-limiti-notevoli-trigonometrici}
permettono di calcolare le derivate delle principali
funzioni elementari.

\begin{proposition}[Derivate delle funzioni fondamentali]
\label{prop:lez22-derivate-fondamentali}
Valgono le seguenti formule:
\begin{enumerate}
  \item per ogni \(c\in\R\),
  \[
    f(x)=c\quad\Longrightarrow\quad f'(x)=0
    \qquad\forall x\in\R;
  \]
  \item per ogni \(n\in\N\setminus\{0\}\),
  \[
    f(x)=x^n\quad\Longrightarrow\quad f'(x)=nx^{n-1}
    \qquad\forall x\in\R;
  \]
  \item per ogni \(\alpha\in\R\),
  \[
    f(x)=x^\alpha\quad\Longrightarrow\quad
    f'(x)=\alpha x^{\alpha-1}
    \qquad\forall x>0;
  \]
  \item per ogni \(a>0\),
  \[
    f(x)=a^x\quad\Longrightarrow\quad
    f'(x)=a^x\log a
    \qquad\forall x\in\R;
  \]
  \item
  \[
    f(x)=e^x\quad\Longrightarrow\quad f'(x)=e^x
    \qquad\forall x\in\R;
  \]
  \item
  \[
    f(x)=\sin x\quad\Longrightarrow\quad
    f'(x)=\cos x
    \qquad\forall x\in\R;
  \]
  \item
  \[
    f(x)=\cos x\quad\Longrightarrow\quad
    f'(x)=-\sin x
    \qquad\forall x\in\R;
  \]
  \item per ogni \(a>0\), con \(a\ne1\),
  \[
    f(x)=\log_a x\quad\Longrightarrow\quad
    f'(x)=\frac{\log_a e}{x}
    \qquad\forall x>0.
  \]
\end{enumerate}
\end{proposition}

\begin{proof}[Dimostrazione del primo punto]
Per \(f(x)=c\), per ogni \(x\in\R\) e \(h\ne0\),
si ha
\[
  \frac{f(x+h)-f(x)}h=\frac{c-c}{h}=0.
\]
Pertanto \(f'(x)=0\) per ogni \(x\in\R\).
\end{proof}

\begin{proof}[Dimostrazione del secondo punto]
\emph{Primo caso: \(n=1\).}
Poiché \(f(x)=x\), per ogni \(h\ne0\),
\[
  \frac{f(x+h)-f(x)}h=\frac{(x+h)-x}{h}=1,
\]
e quindi \(f'(x)=1\).

\emph{Secondo caso: \(n\ge2\).}
Per \(x\ne0\), posto \(t:=h/x\), si ottiene
\[
\begin{aligned}
  f'(x)
  &=\lim_{h\to0}\frac{(x+h)^n-x^n}{h}\\
  &=x^{n-1}\lim_{t\to0}\frac{(1+t)^n-1}{t}
   =nx^{n-1},
\end{aligned}
\]
per la \cref{prop:lez19-limite-potenze}.
Per \(x=0\), invece,
\[
  f'(0)=\lim_{h\to0}\frac{h^n}{h}
  =\lim_{h\to0}h^{n-1}=0.
\]
Pertanto \(f'(x)=nx^{n-1}\) per ogni \(x\in\R\)
e ogni \(n\ge1\).
\end{proof}

\begin{proof}[Dimostrazione del terzo punto]
Siano \(\alpha\in\R\) e \(x>0\). Per \(h\)
sufficientemente piccolo si ha \(x+h>0\). Ponendo
\(t:=h/x\), per la
\cref{prop:lez19-limite-potenze} si ottiene
\[
\begin{aligned}
  f'(x)
  &=\lim_{h\to0}\frac{(x+h)^\alpha-x^\alpha}{h}\\
  &=x^{\alpha-1}
    \lim_{t\to0}\frac{(1+t)^\alpha-1}{t}\\
  &=\alpha x^{\alpha-1}.
\end{aligned}
\]
\end{proof}

\begin{proof}[Dimostrazione del quarto punto]
Siano \(a>0\) e \(x\in\R\). Per ogni \(h\ne0\),
\[
  \frac{a^{x+h}-a^x}{h}
  =a^x\frac{a^h-1}{h}.
\]
Per il quinto punto del
\cref{thm:lez17-limiti-notevoli}, segue
\[
  f'(x)=a^x\lim_{h\to0}\frac{a^h-1}{h}
  =a^x\log a.
\]
Per \(a=1\), il risultato è \(0\), come per
le funzioni costanti.
\end{proof}

\begin{proof}[Dimostrazione del quinto punto]
Ponendo \(a=e\) nel quarto punto e ricordando che
\(\log e=1\), si ha
\[
  \frac{d}{dx}e^x=e^x\log e=e^x.
\]
\end{proof}

\begin{proof}[Dimostrazione del sesto punto]
Per le formule di prostaferesi,
\[
\begin{aligned}
  \frac{\sin(x+h)-\sin x}{h}
  &=\frac{2\sin(h/2)\cos(x+h/2)}{h}\\
  &=\frac{\sin(h/2)}{h/2}\cos\left(x+\frac h2\right).
\end{aligned}
\]
Ponendo \(u:=h/2\), dal
\cref{thm:lez20-limite-seno} e dalla
\cref{prop:lez20-limiti-trigonometrici} segue
\[
  f'(x)=\lim_{u\to0}\frac{\sin u}{u}
    \cos(x+u)=\cos x.
\]
\end{proof}

\begin{proof}[Dimostrazione del settimo punto]
Per le formule di prostaferesi,
\[
\begin{aligned}
  \frac{\cos(x+h)-\cos x}{h}
  &=-\frac{2\sin(x+h/2)\sin(h/2)}{h}\\
  &=-\frac{\sin(h/2)}{h/2}\sin\left(x+\frac h2\right).
\end{aligned}
\]
Ponendo \(u:=h/2\), dal
\cref{thm:lez20-limite-seno} e dalla
\cref{prop:lez20-limiti-trigonometrici} segue
\[
  f'(x)=-\lim_{u\to0}\frac{\sin u}{u}
    \sin(x+u)=-\sin x.
\]
\end{proof}

\begin{proof}[Dimostrazione dell'ottavo punto]
Siano \(a>0\), con \(a\ne1\), e \(x>0\).
Per \(h\) sufficientemente piccolo,
\(1+h/x>0\), e dalle proprietà dei logaritmi segue
\[
\begin{aligned}
  \frac{\log_a(x+h)-\log_a x}{h}
  &=\frac{\log_a(1+h/x)}{h}\\
  &=\frac1x\,\frac{\log_a(1+h/x)}{h/x}.
\end{aligned}
\]
Posto \(t:=h/x\), per il quarto punto del
\cref{thm:lez17-limiti-notevoli}, si ottiene
\[
  f'(x)=\frac1x\lim_{t\to0}
    \frac{\log_a(1+t)}t=\frac{\log_a e}{x}.
\]
\end{proof}
